\section{Trusted Code}
\label{sec:trusted}

Because runtime checks impose a performance overhead, the VM supports \emph{trusted code}: modules for which these checks are waived.
Trusted deployment is currently under Aptos Foundation governance and limited to the Aptos framework; each deployment is approved by the Aptos operators, and we are also investigating staking-based models.
Unlike arbitrary bytecode published on-chain, trusted code is always compiled from source, and thus passes the compiler's extensive type, ability, and reference safety checks before undergoing static bytecode verification.
Defense in depth is therefore preserved: for trusted code, the compiler and the static verifier form the two independent layers, with the compiler playing the role that the runtime checks play for untrusted code.

Executing trusted code disables the checks of \S\ref{sec:type} and \S\ref{sec:ability}.
Because checks like type safety accumulate state during execution, the transitions between untrusted and trusted code must be modeled carefully.
For example, trusted code does not maintain the type stack, so when returning from a trusted function the return value's type must be synthesized for the untrusted caller from the type signature.

%\begin{itemize}
%\item $\CALL trusted$ from untrusted code:
%the types of the arguments need to be removed from $\theta$ since the callee will not use them.
%\item $\RET$ from trusted to untrusted code: the types of the return values need to be pushed on the stack, since %the trusted code does not do this but the untrusted expects it.
%\item $\CALL untrusted$ from trusted code: argument types need to be pushed onto the type stack.
%\item $\RET$ from untrusted code to trusted code: the return types need to be popped.
%\end{itemize}

%%% Local Variables:
%%% mode: latex
%%% TeX-master: "main"
%%% End:
